\documentclass[11pt]{article}

% --- Idioma y fuentes ------------------------------------------------
\usepackage[spanish,es-noshorthands,es-tabla]{babel}
\usepackage[T1]{fontenc}
\usepackage{lmodern}

% --- Geometría y encabezado ------------------------------------------
\usepackage[a4paper,margin=2.5cm]{geometry}
\usepackage{fancyhdr}
\setlength{\headheight}{38.22325pt}
\addtolength{\topmargin}{-6.22325pt}

% --- Matemáticas ------------------------------------------------------
\usepackage{amsmath,amssymb}

% --- Colores, tablas, listas, código ----------------------------------
\usepackage[dvipsnames,table]{xcolor}
\usepackage{booktabs}
\usepackage{array}
\usepackage{longtable}
\usepackage{enumitem}
\usepackage{listings}
\renewcommand{\arraystretch}{1.5}

\lstset{
  language=R,
  basicstyle=\ttfamily\footnotesize,
  commentstyle=\color{gray}\itshape,
  keywordstyle=\color{blue!70!black},
  stringstyle=\color{teal!60!black},
  breaklines=true,
  columns=fullflexible,
  showstringspaces=false
}


\pagestyle{fancy}
\fancyhf{}
\fancyhead[L]{\textbf{Econometría II}\\ Trabajo práctico --- Etapa 1}
\fancyhead[C]{Universidad de Santiago de Chile}
\fancyfoot[C]{\thepage}

\begin{document}
\sloppy

% =====================================================================
%  PORTADA
% =====================================================================
\begin{center}
  {\LARGE \textbf{Etapa 1: construcción de la base de datos}}\\[4pt]
  {\large Econometría II --- Admisión 2026}
\end{center}
\vspace{0.5em}

\vspace{0.8cm}

\begin{center}
\textbf{Integrantes:}\quad
Nombre 1 $\bullet$
Nombre 2 $\bullet$
Nombre 3 $\bullet$
Nombre 4

\vspace{0.4cm}

\textbf{Profesor:}\quad Brian Pfeng Ibacache
\qquad
\textbf{Fecha:}\quad 8 de octubre de 2026
\end{center}

\bigskip

% =====================================================================
\subsection*{1 \; Objetivo}
% =====================================================================

El tema del trabajo es la matrícula en educación superior del proceso de
admisión 2026. En esta etapa construimos un solo data frame con una fila
por estudiante matriculado. Reúne los datos del estudiante, sus notas de
enseñanza media, sus puntajes PAES, los datos de su matrícula y variables
del colegio de egreso. Con esta base definimos la variable de respuesta,
que se estimará con un modelo logit en la Etapa 3.

\medskip\hrule\medskip

% =====================================================================
\subsection*{2 \; Fuentes de datos}
% =====================================================================

Todos los datos vienen del portal de bases de datos del DEMRE, proceso de
admisión 2026. Usamos cuatro archivos, que se unen por el identificador
del estudiante, \texttt{ID\_aux}.

\begin{center}
\begin{tabular}{>{\raggedright\arraybackslash}p{3.4cm}
                >{\raggedright\arraybackslash}p{6.4cm}
                >{\raggedleft\arraybackslash}p{2.6cm}}
\toprule
\textbf{Archivo} & \textbf{Contenido} & \textbf{Filas} \\
\midrule
B (\texttt{df1}) & Inscritos: colegio, domicilio, sexo, ingreso familiar & 317.700 \\
C (\texttt{df2}) & Notas, NEM, ranking y puntajes PAES & 317.700 \\
D (\texttt{df3}) & Postulaciones (una fila por carrera postulada) & 1.686.689 \\
Matrícula (\texttt{df4}) & Vía, tipo de matrícula, preferencia y puntaje & 128.345 \\
\bottomrule
\end{tabular}
\end{center}

No usamos los archivos K y L (cuestionario de caracterización), porque el
DEMRE indica que no se pueden unir con las otras bases.

\medskip\hrule\medskip

% =====================================================================
\subsection*{3 \; Construcción de la base}
% =====================================================================

El código está en R y usa \texttt{tidyverse}. Está completo en el anexo.

\subsubsection*{3.1 \; Variables del colegio}

Antes de filtrar, calculamos las variables del colegio agrupando por
establecimiento (\texttt{RBD}). Se calculan antes porque necesitan a todos
los estudiantes del colegio, no solo a los que quedan en la base final.

\begin{itemize}[leftmargin=1.5em, itemsep=0.3em]
  \item Del archivo B: número de inscritos del colegio y promedio del
        ingreso per cápita familiar de sus inscritos.
  \item Del archivo de matrícula: número de matriculados por vía regular y
        tipo de matrícula 1 que egresaron del colegio. Los colegios sin
        matriculados quedan con 0.
\end{itemize}

\subsubsection*{3.2 \; Selección de columnas y unión}

De cada archivo dejamos solo las columnas que usamos (detalle en la
sección 7). El archivo D trae una fila por carrera postulada, así que lo
dejamos en una fila por estudiante (188.538 estudiantes). El archivo de
matrícula ya trae una fila por estudiante.

Luego unimos los cuatro archivos por \texttt{ID\_aux} con
\texttt{full\_join}, lo que da 317.700 filas, y a esa base le aplicamos
los filtros.

\subsubsection*{3.3 \; Filtros aplicados}

Los filtros que pide el enunciado, y cómo se aplicaron:

\begin{enumerate}[leftmargin=1.5em, itemsep=0.3em, label=\textbf{\alph*)}]
  \item \textbf{Promoción del año:} estudiantes que egresaron de enseñanza
        media en 2025 (\texttt{ANYO\_EGRESO} igual a 2025) y que
        inscribieron al menos una prueba del proceso 2026
        (\texttt{RINDIO\_PROCESO\_ACTUAL} igual a 1).
  \item \textbf{Vía y tipo de matrícula:} matriculados por vía regular
        (\texttt{VIA} igual a 1) con tipo de matrícula «postulación
        regular» (\texttt{TIPO\_MATRICULA} igual a 1).
  \item \textbf{Preferencia:} matriculados en su primera, segunda o
        tercera preferencia (\texttt{PREFERENCIA} entre 1 y 3).
\end{enumerate}

Además, sacamos a quienes no tienen puntaje NEM o lo tienen igual a 0,
porque sin notas de enseñanza media no se pueden comparar con el resto.

\subsubsection*{3.4 \; Observaciones en cada paso}

\begin{center}
\begin{tabular}{>{\raggedright\arraybackslash}p{9cm}
                >{\raggedleft\arraybackslash}p{3cm}}
\toprule
\textbf{Paso} & \textbf{Observaciones} \\
\midrule
Unión de los cuatro archivos & 317.700 \\
Egresados en 2025 & 218.989 \\
Vía regular y tipo de matrícula 1 & 69.004 \\
Inscribieron alguna prueba del proceso 2026 & 69.002 \\
Matriculados en preferencia 1 a 3 & 60.558 \\
Con puntaje NEM distinto de 0 & 60.544 \\
\rowcolor{yellow!30}
\textbf{Base final} & \textbf{60.544} \\
\bottomrule
\end{tabular}
\end{center}

\medskip\hrule\medskip

% =====================================================================
\subsection*{4 \; Variable de respuesta}
% =====================================================================

Por los filtros, todos los estudiantes de la base se matricularon en una
de sus tres primeras preferencias. Por eso la variable de respuesta separa
a quienes quedaron en su primera preferencia de quienes quedaron en la
segunda o la tercera:

\[
\texttt{PREF\_PRIMERA\_DUMMY}_i =
\begin{cases}
1 & \text{si se matriculó en su primera preferencia} \\[4pt]
0 & \text{si se matriculó en su segunda o tercera preferencia}
\end{cases}
\]

\begin{center}
\begin{tabular}{lrr}
\toprule
\textbf{Valor} & \textbf{Estudiantes} & \textbf{Porcentaje} \\
\midrule
1 (primera preferencia) & 40.173 & 66,4\% \\
0 (segunda o tercera)   & 20.371 & 33,6\% \\
\midrule
\textbf{Total} & \textbf{60.544} & 100\% \\
\bottomrule
\end{tabular}
\end{center}

De los 20.371 que no quedaron en su primera preferencia, 13.322 se
matricularon en la segunda y 7.049 en la tercera.

\medskip\hrule\medskip

% =====================================================================
\subsection*{5 \; Variables creadas}
% =====================================================================

\subsubsection*{5.1 \; Variables del establecimiento}

Se pegan a cada estudiante por el \texttt{RBD} de su colegio de egreso.
Así cada estudiante tiene información de su colegio en el mismo proceso
2026.

\begin{itemize}[leftmargin=1.5em, itemsep=0.3em]
  \item \texttt{N\_ALUMNOS}: tamaño del colegio, medido como número de
        inscritos en el proceso que egresaron de él.
  \item \texttt{INGRESO\_COLEGIO}: promedio del tramo de ingreso per
        cápita familiar de los inscritos del colegio.
  \item \texttt{N\_MATRICULADOS}: número de egresados del colegio que se
        matricularon por vía regular.
\end{itemize}

\subsubsection*{5.2 \; Variables dummy}

\begin{itemize}[leftmargin=1.5em, itemsep=0.3em]
  \item \texttt{PREF\_PRIMERA\_DUMMY}: variable de respuesta (sección 4).
  \item \texttt{TP\_DUMMY}: 1 si el colegio es técnico profesional
        (rama T1 a T5), 0 si no. Hay 9.110 estudiantes de colegios
        técnico profesionales y 51.434 del resto.
  \item \texttt{TIPO\_PREF\_DUMMY}: 1 si el tipo de postulación
        registrado es «REGULAR», 0 si es BEA, PACE o género. Vale 1 en
        42.124 estudiantes y 0 en 18.420.
\end{itemize}

\medskip\hrule\medskip

% =====================================================================
\subsection*{6 \; Datos sin información}
% =====================================================================

\begin{center}
\begin{tabular}{>{\raggedright\arraybackslash}p{5.4cm}
                >{\raggedleft\arraybackslash}p{1.8cm}
                >{\raggedright\arraybackslash}p{5.6cm}}
\toprule
\textbf{Variable} & \textbf{Casos} & \textbf{Qué significa} \\
\midrule
\texttt{RBD}, \texttt{RAMA\_EDUCACIONAL}, \texttt{GRUPO\_DEPENDENCIA},
\texttt{CODIGO\_COMUNA} & 244 & Estudiante sin colegio de egreso registrado. \\
\texttt{CODIGO\_REGION} & 10 & Región del colegio sin dato. \\
\texttt{PACE} & 55.382 & No fue preseleccionado PACE. \\
\texttt{BEA} & 53.805 & No fue preseleccionado a la Beca de Excelencia Académica. \\
\texttt{INGRESO\_PERCAPITA\_GRUPO\_FA} igual a 99 & 16.235 & «Prefiero no responder». \\
\bottomrule
\end{tabular}
\end{center}

En el resto de las variables no hay datos faltantes.

\medskip\hrule\medskip

% =====================================================================
\subsection*{7 \; Definición de las variables}
% =====================================================================

Fuente: B, C, D y M (matrícula) son los archivos del DEMRE; «creada»
indica que la construimos nosotros.

\begin{longtable}{>{\raggedright\arraybackslash}p{5.4cm}
                  >{\raggedright\arraybackslash}p{1.4cm}
                  >{\raggedright\arraybackslash}p{6.6cm}}
\toprule
\textbf{Variable} & \textbf{Fuente} & \textbf{Definición} \\
\midrule
\endfirsthead
\toprule
\textbf{Variable} & \textbf{Fuente} & \textbf{Definición} \\
\midrule
\endhead
\bottomrule
\endfoot
\multicolumn{3}{l}{\textit{Identificación y respuesta}} \\
\texttt{ID\_aux} & B, C, D, M & Identificador único del estudiante. \\
\texttt{PREF\_PRIMERA\_DUMMY} & Creada & Variable de respuesta: 1 si se matriculó en su primera preferencia, 0 si en la segunda o tercera. \\
\midrule
\multicolumn{3}{l}{\textit{Estudiante}} \\
\texttt{FECHA\_NACIMIENTO} & B & Mes y año de nacimiento, en formato MMAAAA. \\
\texttt{SEXO} & B & 1 hombre, 2 mujer. \\
\texttt{CODIGO\_COMUNA\_D} & B & Comuna de domicilio del estudiante. \\
\texttt{INGRESO\_PERCAPITA\_GRUPO\_FA} & B & Tramo de ingreso per cápita del grupo familiar, de 1 a 10; 99 es «prefiero no responder». \\
\texttt{PACE} & B & Preseleccionado PACE por el MINEDUC. \\
\texttt{BEA} & B & Preseleccionado a la Beca de Excelencia Académica. \\
\texttt{ANYO\_EGRESO} & B & Año de egreso de enseñanza media (2025 en toda la base). \\
\texttt{RINDIO\_PROCESO\_ACTUAL} & B & 1 si inscribió al menos una prueba del proceso 2026 (1 en toda la base). \\
\midrule
\multicolumn{3}{l}{\textit{Notas y puntajes}} \\
\texttt{PROMEDIO\_NOTAS} & C & Promedio de notas de enseñanza media. \\
\texttt{PTJE\_NEM} & C & Puntaje NEM (notas de enseñanza media). \\
\texttt{PTJE\_RANKING} & C & Puntaje ranking de notas. \\
\texttt{CLEC\_REG\_ACTUAL} & C & Puntaje PAES de Competencia Lectora, rendición regular 2026. \\
\texttt{MATE1\_REG\_ACTUAL} & C & Puntaje PAES de Competencia Matemática 1, rendición regular 2026. \\
\midrule
\multicolumn{3}{l}{\textit{Postulación y matrícula}} \\
\texttt{TIPO\_PREF} & D & Tipo de postulación: REGULAR, BEA, PACE o GÉNERO. \\
\texttt{TIPO\_PREF\_DUMMY} & Creada & 1 si \texttt{TIPO\_PREF} es REGULAR, 0 si no. \\
\texttt{PREFERENCIA} & M & Preferencia en que se matriculó (1 a 3 en la base). Con ella se construye la variable de respuesta. \\
\texttt{VIA} & M & Vía de matrícula (1, regular, en toda la base). \\
\texttt{TIPO\_MATRICULA} & M & Tipo de matrícula (1, postulación regular, en toda la base). \\
\texttt{PTJE\_POND} & M & Puntaje ponderado con que se matriculó. \\
\midrule
\multicolumn{3}{l}{\textit{Colegio de egreso}} \\
\texttt{RBD} & B & Identificador del establecimiento. \\
\texttt{RAMA\_EDUCACIONAL} & B & H1 y H2 humanista científico (diurno y nocturno); T1 a T5 técnico profesional. \\
\texttt{GRUPO\_DEPENDENCIA} & B & 1 particular pagado, 2 particular subvencionado, 3 municipal, 4 servicio local de educación. \\
\texttt{CODIGO\_REGION} & B & Región del establecimiento. \\
\texttt{CODIGO\_COMUNA} & B & Comuna del establecimiento. \\
\texttt{TP\_DUMMY} & Creada & 1 si el colegio es técnico profesional (T1 a T5), 0 si no. \\
\texttt{N\_ALUMNOS} & Creada & Inscritos del proceso que egresaron del colegio. \\
\texttt{INGRESO\_COLEGIO} & Creada & Promedio de \texttt{INGRESO\_PERCAPITA\_GRUPO\_FA} de los inscritos del colegio. \\
\texttt{N\_MATRICULADOS} & Creada & Egresados del colegio matriculados por vía regular y tipo de matrícula 1. \\
\end{longtable}

\texttt{PREFERENCIA} queda en la base como información de la matrícula,
pero no se usará como variable explicativa, porque la variable de
respuesta se construye con ella.

\medskip\hrule\medskip

% =====================================================================
\subsection*{8 \; Base final}
% =====================================================================

La base final (\texttt{df\_final}) tiene \textbf{60.544 filas y 30
columnas}, con una fila por estudiante: cada \texttt{ID\_aux} aparece una
sola vez. Los estudiantes vienen de 3.108 colegios distintos.

\begin{center}
\begin{tabular}{lrrrr}
\toprule
\textbf{Variable} & \textbf{Mínimo} & \textbf{Mediana} & \textbf{Promedio} & \textbf{Máximo} \\
\midrule
\texttt{PROMEDIO\_NOTAS}   & 4,50  & 6,33  & 6,27  & 7,00 \\
\texttt{PTJE\_NEM}         & 259   & 808   & 792   & 1.000 \\
\texttt{PTJE\_RANKING}     & 259   & 822   & 806   & 1.000 \\
\texttt{CLEC\_REG\_ACTUAL} & 316   & 663   & 667   & 1.000 \\
\texttt{MATE1\_REG\_ACTUAL}& 320   & 682   & 698   & 1.000 \\
\texttt{PTJE\_POND}        & 354,6 & 716,1 & 718,6 & 998,4 \\
\texttt{N\_ALUMNOS}        & 8     & 127   & 172   & 3.169 \\
\texttt{INGRESO\_COLEGIO}  & 1,0   & 29,7  & 31,8  & 94,1 \\
\texttt{N\_MATRICULADOS}   & 1     & 55    & 74    & 902 \\
\bottomrule
\end{tabular}
\end{center}

En la base hay 31.981 mujeres (52,8\%) y 28.563 hombres (47,2\%).

% =====================================================================
\newpage
\subsection*{Anexo \; Código en R}
% =====================================================================

\begin{lstlisting}
library(tidyverse)
rm(list = ls())
cat("\014")

df1 <- read_csv2("/Users/david/Downloads/econometria-ii-demre/datos/crudos/base de datos/ArchivoB_Adm2026REG.csv")
df2 <- read_csv2("/Users/david/Downloads/econometria-ii-demre/datos/crudos/base de datos/ArchivoC_Adm2026REG.csv")
df3 <- read_csv2("/Users/david/Downloads/econometria-ii-demre/datos/crudos/base de datos/ArchivoD_Adm2026REG.csv")
df4 <- read_csv2("/Users/david/Downloads/econometria-ii-demre/datos/crudos/base de datos/ArchivoMatr_Adm2026.csv")



############## VARIABLES DEL COLEGIO #######################
colegio <- df1 %>%
  group_by(RBD) %>%
  summarise(N_ALUMNOS = n(),
            INGRESO_COLEGIO = mean(INGRESO_PERCAPITA_GRUPO_FA, na.rm = TRUE))
colegio

matricula<- df4 %>%
  filter(VIA==1 & TIPO_MATRICULA==1) %>%
  left_join(df1 %>% select(ID_aux, RBD), by = "ID_aux") %>%
  group_by(RBD) %>%
  summarise(N_MATRICULADOS = n())
matricula

colegio <- colegio %>%
  left_join(matricula, by = "RBD")
colegio$N_MATRICULADOS <- ifelse(is.na(colegio$N_MATRICULADOS), 0, colegio$N_MATRICULADOS)
colegio



############## FILTRAR DATOS #######################
df1 <- df1 %>%
  select(ID_aux,ANYO_EGRESO,FECHA_NACIMIENTO,RBD,RAMA_EDUCACIONAL,GRUPO_DEPENDENCIA,CODIGO_REGION,CODIGO_COMUNA,CODIGO_COMUNA_D,SEXO,PACE,BEA,INGRESO_PERCAPITA_GRUPO_FA,RINDIO_PROCESO_ACTUAL)
df1

df2 <- df2%>%
  select(ID_aux,PROMEDIO_NOTAS,PTJE_NEM, PTJE_RANKING, CLEC_REG_ACTUAL,MATE1_REG_ACTUAL)
df2

df3 <- df3 %>%
  select(ID_aux, TIPO_PREF) %>%
  distinct(ID_aux, .keep_all = TRUE)
df3

df4 <- df4 %>%
  select(ID_aux, PREFERENCIA,VIA, PTJE_POND,TIPO_MATRICULA) %>%
  distinct(ID_aux, .keep_all = TRUE)
df4


lista_dfs <- list(df1, df2, df3,df4)

df_final <- lista_dfs %>%
  reduce(full_join, by = "ID_aux")%>%
  filter(ANYO_EGRESO==2025)%>%
  filter(VIA==1 & TIPO_MATRICULA==1)%>%
  filter(RINDIO_PROCESO_ACTUAL==1)%>%
  filter(PREFERENCIA %in% c(1,2,3))%>%
  filter(!is.na(PTJE_NEM), PTJE_NEM != 0)
df_final
df_final <- left_join(df_final, colegio, by = "RBD")
df_final

################ CREACION VARIABLES DUMMY ########################

df_final$PREF_PRIMERA_DUMMY <- ifelse(df_final$PREFERENCIA == 1, 1, 0)
df_final$PREF_PRIMERA_DUMMY

df_final$TIPO_PREF_DUMMY<-ifelse(df_final$TIPO_PREF =="REGULAR",1,0)
df_final$TIPO_PREF_DUMMY

df_final$TP_DUMMY <- ifelse(df_final$RAMA_EDUCACIONAL %in% c("T1","T2","T3","T4","T5"), 1, 0)
df_final$TP_DUMMY

table(df_final$PREFERENCIA)
summary(df_final)
\end{lstlisting}

\end{document}
